\documentclass{article}
\usepackage{amsmath}
\usepackage{caption}
\usepackage{graphicx}
\usepackage{fontspec}
\usepackage{url}
\usepackage{xurl}
\setmainfont{SimSun}

\title{系统开发工具基础 第一周实验报告}
\author{涂睿昕---25020007115}
\date{}

\begin{document}
	\maketitle
	
	\section{实验概述}
	本次实验为系统开发工具基础第一周课上实验，实验环境为VirtualBox中的Ubuntu
	虚拟机。
    本次完成四项课上实验：Shell文件处理、Shell管道与文本统计、Git制造并解决合并
	冲突、\LaTeX{}文档编写。
	实验主要练习Linux Shell命令、脚本编写Git 版本控制分支与冲突处理、
	\LaTeX{}文档排版，掌握命令行工具链、版本控制的基本流程。
	
	\section{课上实验题目}
	\subsection{q01：含空格文件名的批量整理}
	\subsubsection{题目要求}
	在 q01 目录下完成目录创建、带空格文件生成、查看绝对路径、复制文件并保留目录结构、
	设置文件与目录权限，生成 inventory.txt 记录文件相对路径与字节大小。
	
	\subsubsection{实现过程与关键命令}
	\begin{verbatim}
		mkdir -p ~/sdt-course/week1/q0{1..4}
		cd ~/sdt-course/week1/q01
		mkdir -p input/docs input/tmp
		printf "alpha\nbeta\n" > "input/docs/notes one.txt"
		echo "hidden" > input/docs/.secret.txt
		touch input/tmp/empty.txt
		printf "2026-08-24 system start\n2026-08-24 init finished\n" > input/run.log
		pwd
		ls -laR input/
		mkdir -p work/25020007115
		cd input
		find . -name "*.txt" -exec cp --parents {} "../work/25020007115/" \;
		cd ..
		ls -laR work/25020007115/
		cd "work/25020007115"
		find . -type d -exec chmod 750 {} \;
		find . -type f -exec chmod 640 {} \;
		cd ../..
		cd "work/25020007115"
		find . -type f -printf "%p %s\n" > inventory.txt
		cat inventory.txt
		cd ../..
	\end{verbatim}
	
	\subsubsection{运行结果}
	\begin{verbatim}
		/home/lin/sdt-course/week1/q01
		
		总计 20
		drwxrwxr-x 4 lin lin 4096  8月 24 16:37 .
		drwxrwxr-x 3 lin lin 4096  8月 24 16:33 ..
		drwxrwxr-x 2 lin lin 4096  8月 24 16:35 docs
		-rw-rw-r-- 1 lin lin   49  8月 24 16:37 run.log
		drwxrwxr-x 2 lin lin 4096  8月 24 16:36 tmp
		
		input/docs:
		总计 16
		drwxrwxr-x 2 lin lin 4096  8月 24 16:35 .
		drwxrwxr-x 4 lin lin 4096  8月 24 16:37 ..
		-rw-rw-r-- 1 lin lin   11  8月 24 16:34 'notes one.txt'
		-rw-rw-r-- 1 lin lin    7  8月 24 16:35 .secret.txt
		
		input/tmp:
		总计 8
		drwxrwxr-x 2 lin lin 4096  8月 24 16:36 .
		drwxrwxr-x 4 lin lin 4096  8月 24 16:37 ..
		-rw-rw-r-- 1 lin lin    0  8月 24 16:36 empty.txt
		
		work/25020007115/:
		docs  tmp
		
		work/25020007115/docs:
		'notes one.txt'
		
		work/25020007115/tmp:
		empty.txt
		
		./inventory.txt 0
		./tmp/empty.txt 0
		./docs/notes one.txt 11
		./docs/.secret.txt 7
	\end{verbatim}

	
	\begin{figure}[htbp]
		\centering
		\includegraphics[width=0.9\textwidth]{q1p1.png}
		\caption{q01运行截图1}
		\label{fig:q01_1}
	\end{figure}
	
	\begin{figure}[htbp]
		\centering
		\includegraphics[width=0.9\textwidth]{q1p2.png}
		\caption{q01运行截图2}
		\label{fig:q01_2}
	\end{figure}
	
	\subsubsection{遇到的问题与解决}
	\begin{itemize}
		\item \textbf{问题1：带空格的文件名操作报错}
		\\ 直接使用命令操作带空格的文件名时，系统会将空格识别为参数分隔符，提示找不到文件。
		\\ \textbf{解决：} 使用双引号包裹包含空格的文件名，或通过 \texttt{find -exec} 传递文件路径，确保空格被正确识别。
		\item \textbf{问题2：初始查看目录时未显示隐藏文件}
		\\ 直接使用 \texttt{ls} 命令无法显示以 \texttt{.} 开头的隐藏文件 \texttt{.secret.txt}。
		\\ \textbf{解决：} 加入 \texttt{-a} 参数，使用 \texttt{ls -la} 查看目录下所有文件（含隐藏文件）。
		\item \textbf{问题3：复制文件时丢失目录结构}
		\\ 直接使用 \texttt{cp} 批量复制时，所有文件被平铺到目标目录，丢失了原有的 \texttt{docs/}、\texttt{tmp/} 层级。
		\\ \textbf{解决：} 给 \texttt{cp} 加上 \texttt{--parents} 参数，保留文件的相对目录结构。
	\end{itemize}
	
	\subsection{q02：随机访问日志统计}
	\subsubsection{题目要求}
	编写脚本 \texttt{analyze.sh}，读取CSV访问日志，统计5xx状态码最多的两条path，输出全部
	数据的平均延迟（保留两位小数）；文件不存在时向标准错误输出报错并返回非零退出码。
	
	\subsubsection{测试数据}
	本次统计使用的\texttt{access.csv}日志数据内容如下：
	\begin{verbatim}
		timestamp,user,path,status,latency_ms
		09:00:01,alice,/api/users,200,120
		09:00:02,bob,/api/orders,500,310
		09:00:03,alice,/api/users,503,180
		09:00:04,carol,/login,500,90
		09:00:05,bob,/api/orders,502,260
		09:00:06,dave,/health,200,20
		09:00:07,alice,/api/orders,500,420
		09:00:08,carol,/api/users,500,150
	\end{verbatim}
	
	\subsubsection{实现过程}
	脚本 \texttt{analyze.sh} 完整代码：
	\begin{verbatim}
		cat > analyze.sh << 'EOF'
		#!/usr/bin/env bash
		
		# 检查参数
		if [ $# -ne 1 ]; then
		echo "Usage: $0 <csv_file_path>" >&2
		exit 1
		fi
		
		FILE="$1"
		
		# 文件不存在时输出错误到stderr，返回非零退出码
		if [ ! -f "$FILE" ]; then
		echo "Error: file '$FILE' not found." >&2
		exit 2
		fi
		
		echo "=== Top 2 paths with most 5xx requests ==="
		# 过滤5xx状态，按path统计次数，降序排列，次数相同按path字典序
		awk -F',' 'NR>1 && $4 ~ /^5[0-9]{2}$/ {cnt[$3]++}
		END {
			for (p in cnt) print cnt[p], p
		}' "$FILE" | sort -k1,1nr -k2,2 | head -n 2
		
		echo ""
		echo "=== Average latency (ms) ==="
		# 跳过表头，计算平均延迟，保留两位小数
		awk -F',' 'NR>1 {sum += $5; n++}
		END {
			if (n > 0) printf "%.2f\n", sum / n
		}' "$FILE"
		EOF
	\end{verbatim}
	
	\subsubsection{运行结果}
	1. 传入正常 \texttt{access.csv} 的输出：
	\begin{verbatim}
		=== Top 2 paths with most 5xx requests ===
		3 /api/orders
		2 /api/users
		
		=== Average latency (ms) ===
		193.75
	\end{verbatim}
	
	2. 传入不存在文件的错误输出：
	\begin{verbatim}
		Error: file 'not_exist.csv' not found.
		Exit code: 2
	\end{verbatim}
	
	\begin{figure}[htbp]
		\centering
		\includegraphics[width=0.9\textwidth]{q2p1.png}
		\caption{q02运行截图1}
		\label{fig:q02_1}
	\end{figure}
	
	\begin{figure}[htbp]
		\centering
		\includegraphics[width=0.9\textwidth]{q2p2.png}
		\caption{q02运行截图2}
		\label{fig:q02_2}
	\end{figure}
	
	\begin{figure}[htbp]
		\centering
		\includegraphics[width=0.9\textwidth]{q2p3.png}
		\caption{q02运行截图3}
		\label{fig:q02_3}
	\end{figure}
	
	\subsubsection{遇到的问题与解决}
	\begin{itemize}
		\item \textbf{问题1：统计结果包含表头行数据}
		\\ 初始编写脚本时未跳过表头，导致第一行的字段名被当作数据参与统计，结果错误。
		\\ \textbf{解决：} 在 \texttt{awk} 中加入 \texttt{NR>1} 条件，跳过第一行表头，只处理数据行。
		\item \textbf{问题2：次数相同时未按路径字典序排序}
		\\ 仅按出现次数降序排序时，次数相同的路径顺序随机，不符合题目要求。
		\\ \textbf{解决：} \texttt{sort} 命令设置多字段排序，先按次数字段降序，再按路径字段升序（字典序）。
		\item \textbf{问题3：错误信息输出到标准输出}
		\\ 初始的错误提示直接用 \texttt{echo} 输出，默认流向标准输出，不符合输出到 stderr 的要求。
		\\ \textbf{解决：} 在 echo 后加上 \texttt{>\&2}，将错误信息重定向到标准错误流。
	\end{itemize}
	
	\subsection{q03：制造、解决并解释一次合并冲突}
	\subsubsection{题目要求}
	初始化git仓库，在main分支创建config.txt；
	创建feature‑a、feature‑b两个分支分别修改配置；
	回到main先合并feature‑a，再合并feature‑b；解决冲突保留 \texttt{mode=safe}，增加一行 \texttt{note=reviewed}，
	最后查看分支提交图。
	
	
	
	\subsubsection{实现过程}
	\begin{verbatim}
		cd ~/sdt-course/week1/q03
		git init
		git branch -M main
		echo "mode=normal" > config.txt
		git add config.txt
		git commit -m "init: add default config"
		git checkout -b feature-a
		echo "mode=safe" > config.txt
		git add config.txt
		git commit -m "feat: set mode to safe"
		git checkout main      
		git checkout -b feature-b
		echo "mode=fast" > config.txt
		git add config.txt
		git commit -m "feat: set mode to fast"
		git checkout main
		git merge feature-a -m "merge feature-a"  # 第一次合并无冲突
		git merge feature-b                       # 第二次合并触发冲突
		git add config.txt
		git commit -m "resolve conflict: keep safe mode, add review note"
		git log --all --graph --decorate --oneline
	\end{verbatim}
	
	\subsubsection{运行结果}
	执行 \texttt{git log --all --graph --decorate --oneline} 的输出：
	\begin{verbatim}
		*  9b8fbf8 (HEAD -> main) resolve conflict: keep safe mode, add review note
		|\
		| * db0a511 (feature-b) feat: set mode to fast
		* | c2b9999 (feature-a) feat: set mode to safe
		|/
		* 9a6f312 init: add default config
	\end{verbatim}
	
	\begin{figure}[htbp]
		\centering
		\includegraphics[width=0.9\textwidth]{q3p1.png}
		\caption{q03运行截图1}
		\label{fig:q03_1}
	\end{figure}
	
	\begin{figure}[htbp]
		\centering
		\includegraphics[width=0.9\textwidth]{q3p2.png}
		\caption{q03运行截图2}
		\label{fig:q03_2}
	\end{figure}
	
	\begin{figure}[htbp]
		\centering
		\includegraphics[width=0.9\textwidth]{q3p3.png}
		\caption{q03运行截图3}
		\label{fig:q03_3}
	\end{figure}
	
	\subsubsection{遇到的问题与解决}
	\begin{itemize}
		\item \textbf{问题1：子目录嵌套Git仓库，造成仓库混乱}
		\\ 曾经在上一级目录误执行\texttt{git init}，q03目录内部嵌套了外层git仓库，执行分支、合并
		命令作用于上层仓库，实验目录内文件不受版本控制。
		\\ \textbf{解决：} 删除错误生成的外层\texttt{.git}文件夹，进入q03目录内部再执行\texttt{git init}，保证仓库
		初始化在实验目录根目录。
		\item \textbf{问题2：快进合并时 -m 参数被忽略，没有生成合并提交}
		\\ 执行\texttt{git merge feature-a -m "merge feature-a"}，Git提示Fast‑forward，指定的
		合并提交信息参数失效。
		\\ \textbf{解决：} 快进合并是Git正常行为，不会新建合并节点，无需额外处理，直接继续后续
		分支操作。
		\item \textbf{问题3：合并产生冲突，编辑冲突文件时没有删除冲突标记}
		\\ 合并\texttt{feature‑b}出现冲突，修改内容后忘记删掉 \texttt{<<<<<<<}、\texttt{=======}、\texttt{>>>>>>>}标记，导
		致配置文件混入git冲突符号。
		\\ \textbf{解决：} 编辑\texttt{config.txt}，完整删除全部冲突分隔标记，保留\texttt{mode=safe}，新增一行
		\texttt{note=reviewed}，保存文件。
		\item \textbf{问题4：解决冲突后忘记add，git仍处于冲突状态}
		\\ 修改完冲突文件直接执行commit，git提示仍有未解决冲突，无法提交。
		\\ \textbf{解决：} 修改完成后先执行\texttt{git add config.txt}把修改加入暂存区，再执行\texttt{git commit}完成合并提交。
		\item \textbf{问题5：无法确认冲突是否真正解决，历史看不清分支走向}
		\\ 完成合并之后，不清楚分支分叉、合并流程是否满足题目要求。
		\\ \textbf{解决：} 使用\texttt{git log --all --graph --decorate --oneline}查看图形化提交历史，核对
		分支节点结构。
	\end{itemize}
	
	\subsection{q04：修复并构建一页技术说明}
	\subsubsection{题目要求}
	\begin{verbatim}
		编写LaTeX文档，添加带label的编号公式并在正文引用；添加2行2列带caption与label的表格并在正文引用；编译生成PDF，保证交叉引用不出现??。
	\end{verbatim}
	
	\subsubsection{实现过程}
	核心代码片段：
	\begin{verbatim}
		\documentclass{article}
		\usepackage{amsmath}
		\usepackage{caption}
		\usepackage{fontspec}
		\setmainfont{WenQuanYi Zen Hei}
		
		\begin{document}
			\title{Tool Report}
			\author{涂睿昕-25020007115}
			\maketitle
			
			\section{Result}
			
			本文采用的累加公式如式\eqref{eq:sum}所示，可用于计算自然数序列的求和结果。
			\begin{equation}
				\label{eq:sum}
				\sum_{i=1}^{n} i = \frac{n(n+1)}{2}
			\end{equation}
			
			实验样本的统计数据如表\ref{tab:sample}所示，包含两组测试的核心指标。
			\begin{table}[htbp]
				\centering
				\caption{实验样本数据}
				\label{tab:sample}
				\begin{tabular}{|c|c|}
					\hline
					样本组 & 数值 \\
					\hline
					Group A & 128 \\
					Group B & 256 \\
					\hline
				\end{tabular}
			\end{table}
			
		\end{document}
	\end{verbatim}
	
	\subsubsection{运行结果}
	编译成功得到 report.pdf，公式、表格正常渲染，交叉引用生效。
	
	\begin{figure}[htbp]
		\centering
		\includegraphics[width=0.9\textwidth]{q4p1.png}
		\caption{q04生成pdf效果截图}
		\label{fig:q04_1}
	\end{figure}
	
	\subsubsection{遇到的问题与解决}
	\begin{itemize}
		\item \textbf{问题1：中文显示为方框乱码}
		\\ 使用默认 pdflatex 引擎编译时，中文字符全部显示为空白方框。
		\\ \textbf{解决：} 引入 \texttt{fontspec} 宏包，指定系统中文字体，切换为 XeLaTeX 引擎编译。
		\item \textbf{问题2：公式和表格引用显示为 ??}
		\\ 第一次编译后，正文里的交叉引用位置显示为 \texttt{??}，没有生成正确编号。
		\\ \textbf{解决：} 执行第二次编译，LaTeX 会读取上一次生成的引用信息，生成正确的编号。
	\end{itemize}
	
	\section{课后练习}
	\begin{enumerate}
		\item 实例1：运行 ls -l / 并观察输出。\\
		命令：\texttt{ls -l /}\\
		输出：
		\begin{verbatim}
			总计 3990620
			lrwxrwxrwx   1 root root          7  4月 22  2024 bin -> usr/bin
			drwxr-xr-x   2 root root       4096  2月 26  2024 bin.usr-is-merged
			drwxr-xr-x   4 root root       4096  8月 30 22:36 boot
			dr-xr-xr-x   2 root root       4096  8月  6  2025 cdrom
			drwxr-xr-x  19 root root       4320  8月 31 20:56 dev
			drwxr-xr-x 145 root root      12288  9月  1 16:02 etc
			drwxr-xr-x   3 root root       4096 10月 22  2025 home
			lrwxrwxrwx   1 root root          7  4月 22  2024 lib -> usr/lib
			lrwxrwxrwx   1 root root          9  4月 22  2024 lib64 -> usr/lib64
			drwxr-xr-x   2 root root       4096  4月  8  2024 lib.usr-is-merged
			drwx------   2 root root      16384 10月 22  2025 lost+found
			drwxr-xr-x   3 root root       4096 10月 23  2025 media
			drwxr-xr-x   2 root root       4096  8月  6  2025 mnt
			drwxr-xr-x   4 root root       4096  8月 27 18:11 opt
			dr-xr-xr-x 295 root root          0  8月 31 20:56 proc
			drwx------   6 root root       4096  8月 27 18:07 root
			drwxr-xr-x  36 root root        960  9月  1 16:02 run
			lrwxrwxrwx   1 root root          8  4月 22  2024 sbin -> usr/sbin
			drwxr-xr-x   2 root root       4096  3月 31  2024 sbin.usr-is-merged
			drwxr-xr-x  17 root root       4096  8月 31 16:16 snap
			drwxr-xr-x   2 root root       4096  8月  6  2025 srv
			-rw-------   1 root root 4086300672 10月 22  2025 swap.img
			dr-xr-xr-x  13 root root          0  8月 31 20:56 sys
			drwxrwxrwt  18 root root       4096  9月  1 21:14 tmp
			drwxr-xr-x  12 root root       4096  8月  6  2025 usr
			drwxr-xr-x  14 root root       4096 10月 22  2025 var
		\end{verbatim}
		
		\item 实例2：写一条命令，输出一个同时包含字面量\$ 、! 和换行符的字符串。\\
		命令：\texttt{echo \$'dollar sign: \$ \textbackslash nexclamation mark: !'
			dollar sign: \$ }\\
		输出：\texttt{dollar sign: \$ 
			exclamation mark: !}
		
        \item 实例3：\$? 保存上一条命令的退出状态（0 表示成功）。\&\& 仅在前一条成功时执行后一条；|| 仅在前一条失败时执行后一条。写一个一行命令：仅当 /tmp/mydir 不存在时才创建它。\\
    	命令：\texttt{[ ! -d "/tmp/mydir" ] \&\& mkdir /tmp/mydir}\\
	    输出：\texttt{无}
	    
	    \item 实例4：使用管道找出你「home 目录」中最常见的 5 种文件扩展名。\\
	    命令：\texttt{find ~ -maxdepth 2 -type f | sed -E 's/.*\textbackslash.([\textasciicircum.]+)\textbackslash\$/\textbackslash1/' | grep -v '/' | sort | uniq -c | sort -nr | head -5}\\
	    输出：
	    \begin{verbatim}
	    	8 c
	    	7 pid
	    	6 go
	    	4 utf8
	    	3 py
	    \end{verbatim}
	    
	    \item 实例5：xargs 会把 stdin 的每一行转换为命令参数。结合 find 和 xargs（不要用 find -exec），找出目录中所有 .sh 文件，并用 wc -l 统计每个文件行数。\\
	    命令：\texttt{find . -type f -name "*.sh" -print0 | xargs -0 wc -l}\\
	    输出：
	    \begin{verbatim}
	    	23 ./.local/share/Trash/files/analyze.sh
	    	8 ./sdt-course/week2/q05/worker.sh
	    	30 ./sdt-course/week1/q02/analyze.sh
	    	14 ./.vscode/extensions/ms-python.debugpy-2026.6.0-linux-x64/bundled/libs
	    	/debugpy/_vendored/pydevd/pydevd_attach_to_process/linux_and_mac/compile_mac.sh
	    	11 ./.vscode/extensions/ms-python.debugpy-2026.6.0-linux-x64/bundled/libs
	    	/debugpy/_vendored/pydevd/pydevd_attach_to_process/linux_and_mac/compile_linux.sh
	    	86 总计
	    \end{verbatim}
	    
	    \item 实例6：用 curl 获取示例数据 https://microsoftedge.github.io/Demos/json-dummy-data/64KB.json，再用 jq 提取 version 大于 6 的人员姓名。\\
	    命令：
	    \begin{verbatim}
	    	sudo apt install curl
	    	curl -s https://microsoftedge.github.io/Demos/json-dummy-data/64KB.json | jq
	    	 '.[] | select(.version > 6) | .name'
	    \end{verbatim}
	    输出：\begin{verbatim}
	    	"Adeel Solangi"
	    	"Aamir Solangi"
	    	"Adil Eli"
	    	"Adil Abro"
	    	"Antía Sixirei"
	    	"Aygul Mutellip"
	    	"Celtia Anes"
	    	"George Mifsud"
	    	"Dialè Meso"
	    	"Doreen Bartolo"
	    	"Ali Ayaz"
	    	"Guzelnur Polat"
	    	"John Falzon"
	    	"Marcos Amboade"
	    	"Shafqat Memon"
	    	"Meladi Papo"
	    	"Sabela Veloso"
	    	"Madule Ledimo"
	    	"Philip Camilleri"
	    	"Hunadi Makgatho"
	    	"Charmaine Abela"
	    	"Tumelò Letamo"
	    	"Tegra Núnez"
	    	"Dilnur Qeyser"
	    	"Iago Peirallo"
	    	"Josephine Balzan"
	    	"Mmathabò Mojapelo"
	    	"Kgabo Lerumo"
	    	"Lawrence Scicluna"
	    	"Joseph Grech"
	    	"Lesetja Theko"
	    	"Martiño Arxíz"
	    	"Malehumò Ledwaba"
	    	"Lajwanti Kumari"
	    	"Maria Sammut"
	    	"Matome Molamo"
	    	"Noa Ervello"
	    	"sayama Amir"
	    	"Mariña Quintá"
	    	"Memet Tursun"
	    	"Rashid Rajput"
	    	"Nicholas Micallef"
	    	"Peter Zammit"
	    	"Maname Mohlare"
	    	"Tshepè Mobu"
	    	"Monica Lohana"
	    	"Joanne Scerri"
	    	"Ratanang Maphutha"
	    	"Thobile Mbele"
	    	"Kristján Kristjánsson"
	    	"Stefán Stefánsson"
	    	"Preeti Rajdan"
	    	"Sanjay Trivedi"
	    	"Damir Benic"
	    	"Sigrún Kristjánsdóttir"
	    	"Rajesh Santoshi"
	    	"Mxolisi Mhlongo"
	    	"Monty Dubey"
	    	"Richa Choukse"
	    	"Ingibjörg Ólafsdóttir"
	    	"Mogorosi Bakwena"
	    	"Ronak Gupta"
	    	"Chetana Hegde"
	    	"Mohan Pandey"
	    	"Haris Osmanovic"
	    	"Kenosi Kwenaemang"
	    	"Adnan Spahic"
	    	"Kabelo Morwe"
	    	"Einar Einarsson"
	    	"Sigríður Einarsdóttir"
	    	"Sonu Jain"
	    	"Boitumelo Ngwako"
	    	"Modise Tau"
	    	"Gunnar Gunnarsson"
	    	"Kgosietsile Bogatsu"
	    	"Monika Nayak"
	    	"Lucky Shastry"
	    	"Raju Rathore"
	    	"Meenakshi Benjaree"
	    	"Samir Simic"
	    	"Prashant Chourey"
	    	"Prakash Malviya"
	    	"Ivana Kalic"
	    	"Denis Terzic"
	    	"Bukhosi Bhengu"
	    	"Siyabonga Sithole"
	    	"Rohini Vasav"
	    	"Sunil Kapoor"
	    	"Zamokuhle Zulu"
	    \end{verbatim}
	    
	    \item 实例7：请写一个 awk 命令：只输出第二列大于 100 的行，并交换第一列和第三列。\\
	    命令：\texttt{printf 'a 50 x\textbackslash nb 150 y\textbackslash nc 200 z\textbackslash n' | awk '\$2>100 {print \$3,\$2,\$1}'}\\
	    输出：\texttt{y 150 b
	    	z 200 c}
	    
	    \item 实例8：写一条命令，把文件复制为带当天日期的备份文件名。\\
	    命令：\texttt{cp test.txt test\_\$(date +\%Y-\%m-\%d).txt}\\
	    输出：\texttt{生成test\_2026-09-01.txt}
	    
	    \item 实例9：在命令 find ~/Downloads -type f -name "*.zip" -mtime +30 中，*.zip 是一个 「glob」。什么是 glob ？新建一个测试目录并创建一些文件，试试 ls *.txt 、ls file?.txt 、ls {a,b,c}.txt 等模式。\\
	    命令：
	    \begin{verbatim}
	        mkdir -p ~/globtest && cd ~/globtest
	        touch file1.txt file2.txt filea.txt fileb.txt filec.txt
	        ls *.txt
	        ls file?.txt
	        ls file{a,b,c}.txt
    	\end{verbatim}
	    输出：
	    \begin{verbatim}
	        file1.txt  file2.txt  filea.txt  fileb.txt  filec.txt
	        file1.txt  file2.txt  filea.txt  fileb.txt  filec.txt
	        filea.txt  fileb.txt  filec.txt
	    \end{verbatim}
	    
	    \item 实例10：确认当前 Shell 是否合适，可运行 echo \$SHELL；若输出类似 /bin/bash 或 /usr/bin/zsh ，就说明没问题。\\
	    命令：\texttt{echo \$SHELL}\\
	    输出：\texttt{/bin/bash}
	\end{enumerate}
	
	\section{解题感悟}
	通过第一周的实验，我熟悉了Ubuntu Shell基础命令，学会处理
	含空格文件名、管道工具 \texttt{awk sort head} 等文本处理工具，能够编写
	简单shell脚本完成统计任务。
	
	在Git版本控制部分，理解了分支创建、切换、合并操作，实际体验了合并冲突的产生
	过程以及手动解决冲突的流程，认识到版本控制不应当一次性大批量提交代码。
		
	在LaTeX部分，体会到不同编译引擎之间的差异，pdflatex对中文支持有限，xelatex需要配置
	中文字体。宏包缺失、字体缺失、拼写错误都会直接导致编译失败，排错需要仔细阅读日志输出。
		
	实验过程也体会到课程强调的“可复现、可协作”，所有操作尽量使用命令行完成，每完成
	一道题目进行一次git提交，保留完整的版本历史。
	
	
	\section{代码仓库提交记录}
	GitHub仓库链接：\url{https://github.com/imnotab/sdtb/tree/main}
	
	提交历史截图：
	\begin{figure}[htbp]
		\centering
		\includegraphics[width=0.9\textwidth]{week1commit.png}
		\caption{项目commit历史记录}
	\end{figure}
	
\end{document}
